\documentclass[11pt]{article}
\usepackage[a4paper,margin=1in]{geometry}
\usepackage{amsmath,amssymb,amsthm,booktabs,url}
\theoremstyle{definition}
\newtheorem{theorem}{Theorem}[section]
\newtheorem{lemma}[theorem]{Lemma}
\newtheorem{proposition}[theorem]{Proposition}
\newtheorem{definition}[theorem]{Definition}
\newtheorem{remark}[theorem]{Remark}
\newcommand{\gt}{\gamma_t}
\newcommand{\Zinf}{\mathbb{Z}_{\geq0}\cup\{+\infty\}}
\newcommand{\sq}{\mathbin{\square}}
\font\hangulfont="Noto Sans KR"

\title{Exact Total Domination Numbers of Cylindrical Grids of Widths Five, Six, and Seven}
\author{{\hangulfont 박성현}}
\date{}

\begin{document}
\maketitle

\begin{abstract}
We determine the total domination number of the cylindrical grids
\(P_m\sq C_n\) for every circumference \(n\geq3\) when
\(m\in\{5,6,7\}\). The proof represents a total dominating set by a closed
walk in a finite column automaton whose states record selected rows and the
rows still requiring domination from the next column. The graph parameter is
thereby the minimum diagonal entry of a min-plus matrix power. For widths
five, six, and seven we verify exact entrywise matrix identities with periods
four, fourteen, and four, respectively. Exact finite certificates supplied
with the paper establish these identities and the finite prefix values needed
to complete the computer-assisted part of the proof. The resulting formulas
include the exceptional circumferences \(n=12\) for width six and \(n=7,14\)
for width seven.
\end{abstract}

\noindent\textbf{Keywords:} total domination; cylindrical grid; Cartesian product;
finite-state method; min-plus algebra; computer-assisted proof.

\section{Introduction}

Let \(G\) be a graph without isolated vertices. A set \(D\subseteq V(G)\) is
\emph{totally dominating} if every vertex of \(G\), including every vertex
of \(D\), has a neighbor in \(D\). The minimum cardinality of such a set is
the total domination number \(\gt(G)\). In this paper we study the
cylindrical grids \(P_m\sq C_n\), the Cartesian product of a path and a
cycle.

Total domination in grid-like products has a substantial literature. Gravier
studied total domination in products of paths and cycles and related tiling
formulations \cite{Gravier2002}. Exact cylindrical results for smaller fixed
path widths were obtained by Eakawinrujee for widths two, three, and four
\cite{Eakawinrujee2022}; the author's dissertation gives additional
provenance for that cylinder program \cite{EakawinrujeeThesis2023}.
Rectangular grids, ordinary domination, and several
variants provide nearby but distinct problems; for example, the fixed-width
results of Crevals and Ostergard concern \(P_m\sq P_n\), not cylinders
\cite{CrevalsOstergard2017}, while Nandi, Parui, and Adhikari study ordinary
domination \cite{NandiParuiAdhikari2011}.

The results below concern widths not covered by the fixed-width cylindrical
formulas cited above. Hu, Sohn, and Chen study related bundle and toroidal
cases, including $C_m$ bundles over $C_n$ and an exact result for
$C_n\square C_5$; these are important context but are not silently identified
with the cylinder $P_5\sq C_n$ \cite{HuSohnChen2016}. Eakawinrujee's public
dissertation records the exact subfamily
\[
p\equiv1\pmod2,\quad q\equiv0\pmod4
\quad\Longrightarrow\quad
\gt(P_p\sq C_q)=\frac{(p+1)q}{4},
\]
which overlaps the present width-five and width-seven formulas on infinitely
many circumferences \cite{EakawinrujeeThesis2023}. We therefore make the
narrow claim of complete classifications for widths five, six, and seven
across all $n\ge3$ in the audited record, extending rather than erasing these
earlier exact subfamilies. The cyclic boundary condition couples the first and
last columns, and a selected vertex cannot dominate itself. Our approach
handles both issues with a finite-state representation; the semantic reduction
is proved in the paper and the finite identities are supplied as exact
computer-assisted certificates.

The methodological distinction from recent nearby work is also important.
Garzón, Martínez, Moreno, and Puertas already use min-plus algorithms for
2-domination of small-cycle cylinders \cite{GarzonMartinezMorenoPuertas2021}.
Martínez, Castaño-Fernández, and Puertas later use tropical matrix products
for broader 2-domination cylinder bounds
\cite{MartinezCastanoPuertas2024}; that invariant requires two selected
neighbors for vertices outside the set, whereas ordinary total domination
requires one open neighbor for every vertex, including selected vertices. The
shared min-plus language is methodological context, not an equivalence of
parameters.

\subsection{Prior results and scope of the contribution}

The relevant boundary of the literature can be summarized as follows. The
first two rows concern the same total-domination parameter, while the latter
two provide context or a nearby invariant.

\begin{center}
\small
\begin{tabular}{p{0.25\textwidth}p{0.63\textwidth}}
\toprule
source & scope relevant to this paper\\
\midrule
Eakawinrujee (2022) & exact cylinder results for smaller fixed path widths; bounds and related results beyond that range\\
Eakawinrujee dissertation & exact odd path-width subfamily \(q\equiv0\pmod4\), including infinitely many width-five and width-seven cases\\
Hu--Sohn--Chen (2016) & related bundle/toroidal total and paired domination results, including an exact \(C_n\square C_5\) case\\
Garzón--Martínez--Moreno--Puertas (2021) & min-plus algorithms for 2-domination of small-cycle cylinders, a different invariant\\
Martínez--Castaño-Fernández--Puertas (2024) & finite-state/tropical methods for 2-domination of cylinders, a different invariant\\
\bottomrule
\end{tabular}
\normalsize
\end{center}

The present result is not a claim that any one of these ingredients is new
in isolation. Its bounded claim is that the exact values for every
\(n\ge3\) are supplied for each of the three widths \(m=5,6,7\), with the
earlier exact subfamilies and related parameters explicitly separated from
the all-circumference statement.

Our main results are the following.

\begin{theorem}[Width five]\label{thm:w5}
For every \(n\geq3\),
\[
\gt(P_5\sq C_n)=
\begin{cases}
\left\lceil 3n/2\right\rceil+1,&n\equiv2\pmod4,\\
\left\lceil 3n/2\right\rceil,&\text{otherwise}.
\end{cases}
\]
\end{theorem}

\begin{theorem}[Width six]\label{thm:w6}
Let \(\varepsilon_r\), for \(r\in\{0,\ldots,13\}\), be given by
\[
\varepsilon_r=\begin{cases}
0,&r\in\{0,3,5,7,9,10\},\\
1,&r\in\{1,4,6,11,13\},\\
2,&r\in\{2,8\},\\
3,&r=12.
\end{cases}
\]
Then
\[
\gt(P_6\sq C_n)=\left\lceil\frac{12n}{7}\right\rceil
+\varepsilon_{n\bmod14}
\]
for every \(n\geq3\), except that \(\gt(P_6\sq C_{12})=22\).
\end{theorem}

\begin{theorem}[Width seven]\label{thm:w7}
For every \(n\geq3\),
\[
\gt(P_7\sq C_n)=2n,
\]
except \(\gt(P_7\sq C_7)=15\) and \(\gt(P_7\sq C_{14})=30\).
\end{theorem}

The proof has two layers. First, we prove the graph-theoretic equivalence
between total dominating sets and weighted closed walks. Second, a finite
certificate establishes the matrix identities and prefix values. The
thresholded recurrence lemma then shows exactly how those finite facts imply
each all-circumference formula.

For orientation, exact small-circumference cylinder cases are already part of
the prior record discussed by Hu--Sohn--Chen and reproduced in the public
Eakawinrujee dissertation. For the three target widths, those prior
\(C_3\)- and \(C_4\)-values are \(5,6,6,8,6,8\), respectively. The same
values are independently present in our finite prefix, but are not claimed
as new. The table separates these checked small circumferences from the
all-circumference classification supplied below; it does not attribute the
cylinder cases to Hu--Sohn--Chen's separate toroidal \(C_n\square C_5\)
result:
\[
\begin{array}{c|cc}
 m & \gt(P_m\sq C_3) & \gt(P_m\sq C_4)\\ \hline
 5&5&6\\ 6&6&8\\ 7&6&8
\end{array}
\]

\section{Preliminaries}

We use \([m]=\{0,1,\ldots,m-1\}\) for the rows of \(P_m\), and index the
cycle columns by \(\mathbb{Z}/n\mathbb{Z}\). For \(S\subseteq[m]\), define its
open path neighborhood by
\[
N_P(S)=\{i\in[m]:i-1\in S\text{ or }i+1\in S\},
\]
where indices outside \([m]\) are omitted. Thus a row selected in a column
does not automatically dominate the vertex in that same row and column.

For a set \(D\subseteq V(P_m\sq C_n)\), write
\[
S_j=\{i\in[m]:(i,j)\in D\}.
\]
The horizontal neighbors of \((i,j)\) are \((i,j-1)\) and \((i,j+1)\),
and the vertical neighbors are the vertices in rows adjacent to \(i\) in
the same column.

\section{The finite-state model}

For a cyclic sequence of selected row sets \((S_j)\), define the pending set
\[
R_j=[m]\setminus\bigl(S_{j-1}\cup N_P(S_j)\bigr).
\]
The rows in \(R_j\) are precisely those for which the left horizontal and
vertical domination alternatives are absent in column \(j\). Define the fixed
admissible state set
\[
Q_m=\{(S,R):S\subseteq[m],\ R\subseteq[m]\setminus N_P(S),\
\text{and }R=[m]\setminus(U\cup N_P(S))\text{ for some }U\subseteq[m]\}.
\]
Thus the predecessor mask is quantified in the definition but is not part of
the state. Given a state \(q=(S,R)\in Q_m\) and an input mask
\(T\subseteq[m]\), the transition is legal precisely when
\[
R\subseteq T,
\]
and its head pending set is forced to be
\[
R'=[m]\setminus\bigl(S\cup N_P(T)\bigr).
\]
The edge is \((S,R)\to(T,R')\), where the displayed formula defines \(R'\),
provided \(R\subseteq T\) and \((T,R')\in Q_m\); its weight is \(|T|\).
Every edge head is therefore a member of the fixed finite set \(Q_m\).

\begin{theorem}[Finite-state representation]\label{thm:bijection}
For \(m\in\{5,6,7\}\) and \(n\geq3\), labeled total dominating sets of
\(P_m\sq C_n\) are in weight-preserving bijection with labeled length-\(n\)
closed walks in the state graph above. The weight of a walk is the sum of its
transition weights.
\end{theorem}

\begin{proof}
We first prove soundness. Let a cyclic legal walk be given and define
\(D=\{(i,j):i\in S_j\}\). Consider a vertex \((i,j)\). If
\(i\notin R_j\), then by the definition of \(R_j\), either
\(i\in S_{j-1}\) or \(i\in N_P(S_j)\). Hence \((i,j)\) has a selected left
or vertical open neighbor. If \(i\in R_j\), transition legality gives
\(i\in S_{j+1}\), so the right horizontal neighbor is selected. These cases
include selected vertices themselves and therefore establish total
domination.

Conversely, let \(D\) be a total dominating set and define \(S_j\) from its
selected vertices and \(R_j\) by the displayed formula. If \(i\in R_j\),
then the left and every vertical selected-neighbor alternative is absent.
Total domination forces the right neighbor \((i,j+1)\) to be selected, so
\(R_j\subseteq S_{j+1}\). The next pending formula follows from its
definition. The same argument applies to the wraparound transition.

The two constructions are inverse because the pending masks are forced by
the selected masks. Finally,
\[
\sum_j|S_{j+1}|=\sum_j|S_j|=|D|,
\]
so weights agree. For \(n=3\), the left and right horizontal neighbors are
distinct because \(3\nmid2\), so the cyclic encoding has exactly the simple
graph neighborhood semantics.
\end{proof}

\subsection{A local transition example}

It is useful to see the pending-mask convention in a concrete transition.
Take \(S=\varnothing\) and \(R=[m]\). This is an admissible state: before
the current column is read, every row may still require its right horizontal
neighbor. Choose \(T=[m]\). The edge is legal because \(R\subseteq T\), and
its head pending set is
\[
[m]\setminus\bigl(\varnothing\cup N_P([m])\bigr)=\varnothing.
\]
The transition has weight \(m\). In the corresponding two-column fragment,
the selected vertices in \(T\) are dominated vertically, while the rows that
were pending in the tail receive their required right neighbor. This example
also illustrates why the state records obligations created by the preceding
column rather than treating a selected vertex as self-dominating.

More generally, the state condition is not an additional graph restriction.
For \(R\subseteq[m]\setminus N_P(S)\), taking
\[
U=[m]\setminus\bigl(R\cup N_P(S)\bigr)
\]
gives \(R=[m]\setminus(U\cup N_P(S))\). Thus the quantified predecessor
mask in the definition of \(Q_m\) is exactly the realizability condition for
the residual obligation mask. The independent state reconstruction uses this
equivalent characterization.

\section{Min-plus path semantics and the recurrence certificate}

Let \(M_m\) be the weighted adjacency matrix of this graph over \(\Zinf\),
with \(+\infty\) for a missing edge. Min-plus addition
and multiplication are
\[
a\oplus b=\min(a,b),\qquad a\otimes b=a+b,
\]
with \(+\infty\) absorbing for \(\otimes\). Matrix powers use min-plus
multiplication. The min-plus identity matrix has zero diagonal and
\(+\infty\) off the diagonal. For \(c\in\mathbb Z_{\geq0}\), scalar matrix
shift means \((c\otimes A)_{uv}=c+A_{uv}\) for finite \(A_{uv}\), and
\((c\otimes A)_{uv}=+\infty\) otherwise.

\begin{lemma}[Min-plus path semantics]\label{lem:path}
For every \(k\geq0\), \((M_m^k)_{uv}\) is the minimum weight of a length-\(k\)
walk from \(u\) to \(v\), with the minimum of an empty set interpreted as
\(+\infty\).
\end{lemma}

\begin{proof}
The case \(k=0\) is the min-plus identity matrix. For the induction step,
partition a length-\(k+1\) walk by its penultimate state \(x\). The prefix
has minimum cost \((M_m^k)_{ux}\), the final edge has cost \((M_m)_{xv}\),
and minimizing over the finite state set gives the min-plus product. Missing
prefixes or edges contribute \(+\infty\) and therefore do not create a walk.
\end{proof}

Define
\[
g_m(n)=\min_{q\in Q_m}(M_m^n)_{qq}.
\]
Theorem~\ref{thm:bijection} and Lemma~\ref{lem:path} give
\[
\gt(P_m\sq C_n)=g_m(n),\qquad n\geq3.
\]

The distinction between an entrywise identity and a scalar sequence pattern
is important. An identity for the minimum diagonal alone would not control
arbitrary endpoints and would not justify concatenating an additional walk.
The certificates instead compare every matrix entry at the two exponents.
This is the reason the recurrence lemma applies uniformly to all closed walks,
including those whose minimizing state changes with the circumference.

\begin{lemma}[Thresholded min-plus recurrence]\label{lem:rec}
Suppose an entrywise identity
\[
M^{N+p}=c\otimes M^N
\]
holds for a finite integer \(c\). Define
\(h(t)=\min_q(M^t)_{qq}\). Then for every \(k\geq0\),
\[
M^{N+p+k}=c\otimes M^{N+k},
\]
and hence \(h(n+p)=h(n)+c\) for every finite-value base \(n\geq N\).
\end{lemma}

\begin{proof}
Min-plus matrix multiplication is associative: both parenthesizations
minimize the same finite collection of triple products. A finite scalar shift
commutes with a subsequent min-plus product, including entries equal to
\(+\infty\). Therefore, for \(k\geq0\),
\[
M^{N+p+k}=M^{N+p}\otimes M^k
=(c\otimes M^N)\otimes M^k
=c\otimes(M^N\otimes M^k)
=c\otimes M^{N+k}.
\]
Taking the minimum over diagonal entries gives the same statement for \(g\).
When the relevant diagonal minimum is finite, the scalar operation is ordinary
addition. The conclusion is explicitly restricted to \(n=N+k\geq N\); no
value below the threshold is propagated by this lemma.
\end{proof}

\section{Finite certificates and exact formulas}

The finite certificate proposition below is the computational input to the
three theorem proofs. Its verifier reconstructs the state and transition
semantics, parses the persisted matrices, recomputes both powers, and checks
the entrywise identities with explicit \(+\infty\) semantics. The proposition
is computer-assisted: its trusted inputs are the definitions, the certificate
files, and the verifier specified in the reproducibility record.

\begin{proposition}[Finite certificate and prefix data]\label{prop:finite}
For widths five, six, and seven, the state and transition graphs have the
counts in Table~\ref{tab:cert}; the displayed matrix identities hold; and the
diagonal minima \(g_m(n)\) agree with the exact finite-prefix tables in the
project artifact on \(3\le n\le19\), \(3\le n\le34\), and \(3\le n\le31\),
respectively.
\end{proposition}

\begin{table}[ht]
\centering
\caption{Finite state graphs and entrywise recurrence certificates.}
\label{tab:cert}
\begin{tabular}{rrrrrr}
\toprule
width & states & transitions & \(N\) & \(p\) & \(c\)\\
\midrule
5 & 169 & 2,419 & 16 & 4 & 6\\
6 & 441 & 11,025 & 21 & 14 & 24\\
7 & 1,156 & 50,303 & 28 & 4 & 8\\
\bottomrule
\end{tabular}
\end{table}

The certified identities are
\[
M_5^{20}=6\otimes M_5^{16},\qquad
M_6^{35}=24\otimes M_6^{21},\qquad
M_7^{32}=8\otimes M_7^{28}.
\]
Consequently the recurrences begin at \(n=16\), \(n=21\), and \(n=28\),
respectively. The exact finite-prefix replay covers \(3\leq n\leq19\) for
width five, \(3\leq n\leq34\) for width six, and \(3\leq n\leq31\) for
width seven. Since each interval ends at \(N+p-1\), every later index
reduces to a checked base in \([N,N+p-1]\).

The decisive finite values are displayed here. The complete machine-readable
tables, including all rows from \(3\) through the stated prefix endpoint,
are part of the reproducibility record.

\begin{table}[ht]
\centering
\caption{Width-five prefix values.}\label{tab:w5prefix}
\begin{tabular}{c|rrrrrrrrrrrrrrrrr}
\toprule
\(n\)&3&4&5&6&7&8&9&10&11&12&13&14&15&16&17&18&19\\
\midrule
\(g_5(n)\)&5&6&8&10&11&12&14&16&17&18&20&22&23&24&26&28&29\\
\bottomrule
\end{tabular}
\end{table}

\begin{table}[ht]
\centering
\caption{Width-six base and exceptional values.}\label{tab:w6prefix}
\begin{tabular}{c|rrrrrrrrrrrrrrr}
\toprule
\(n\)&12&21&22&23&24&25&26&27&28&29&30&31&32&33&34\\
\midrule
\(g_6(n)\)&22&36&40&40&42&44&48&48&48&51&54&54&56&57&60\\
\bottomrule
\end{tabular}
\end{table}

\begin{table}[ht]
\centering
\caption{Width-seven exceptional and tail-base values.}\label{tab:w7prefix}
\begin{tabular}{c|rrrrrr}
\toprule
\(n\)&7&14&28&29&30&31\\
\midrule
\(g_7(n)\)&15&30&56&58&60&62\\
\bottomrule
\end{tabular}
\end{table}

\subsection{Width five}

The identity \(M_5^{20}=6\otimes M_5^{16}\) gives
\(g_5(n+4)=g_5(n)+6\) for \(n\geq16\). \textit{Proof of
Theorem~\ref{thm:w5}.} The finite proposition gives every value
\(3\le n\le19\); the base values \(16\le n\le19\) are shown in
Table~\ref{tab:w5prefix}, and the complete prefix comparison is in the
machine-readable table. Every \(n\ge20\) reduces by four to one of these
bases, and the residue \(n\equiv2\pmod4\) contributes the additional unit.

\subsection{Width six}

The identity \(M_6^{35}=24\otimes M_6^{21}\) gives
\(g_6(n+14)=g_6(n)+24\) only for \(n\geq21\). \textit{Proof of
Theorem~\ref{thm:w6}.} The value at \(n=12\) is a finite-prefix exception
and is not propagated to \(n=26\). The finite proposition gives every value
\(3\le n\le34\), including all rows \(3\le n<21\) and every base value
\(21\le n\le34\). For every
\(n\geq35\), repeated subtraction of fourteen reaches one of the checked
values \(21,\ldots,34\), which proves Theorem~\ref{thm:w6}; the recurrence
is never applied below its threshold.

\subsection{Width seven}

The identity \(M_7^{32}=8\otimes M_7^{28}\) gives
\(g_7(n+4)=g_7(n)+8\) only for \(n\geq28\). \textit{Proof of
Theorem~\ref{thm:w7}.} The finite proposition explicitly covers every row
\(3\le n\le31\), including the ordinary prefix rows and the exceptions
\(g_7(7)=15\) and \(g_7(14)=30\). Every \(n\geq32\) reduces to a base in
\(28,\ldots,31\), proving the theorem. The exceptions remain only in the
finite prefix and are not propagated by the recurrence.

\section{Verification and reproducibility}

The repository includes a reference implementation and a separately written
certificate verifier. The verifier rebuilds states, transitions, weights, and
matrix powers from the definitions. Detailed test counts, mutation results,
certificate manifests, and hash bindings are recorded separately from the
mathematical proof.
The three certificate commands are:
\begin{verbatim}
python verify/verify_certificates.py 5 certificates/width5
python verify/verify_certificates.py 6 certificates/width6
python verify/verify_certificates.py 7 certificates/width7
\end{verbatim}
Each command reports zero identity mismatches. The repository’s SHA256SUMS
file binds the finite artifacts, and paper/REPRODUCIBILITY.md records the
expected dimensions, counts, parameters, and verification outputs.

The computations establish the finite proposition; the preceding semantic
lemmas and theorem arguments establish why that proposition implies the
infinite statements.

\subsection{What is and is not certified}

The finite proposition has three logically separate inputs. First, the
definition-level state and transition reconstruction fixes the finite graph.
Second, the persisted matrices are parsed with explicit finite and
\(+\infty\) tags, and the verifier recomputes both powers from the reconstructed
edges rather than trusting a producer's claimed identity. Third, the prefix
trace initializes every state as a possible source and minimizes the matching
diagonal entry, so its values are the unrestricted \(g_m(n)\), not a
low-weight-source surrogate. Hashes bind the persisted binary artifacts to
the manifests used by the verifier.

This separation also fixes the boundary of the computer-assisted argument.
The code supplies finite exact premises; it does not replace the
total-domination bijection, min-plus path lemma, or thresholded algebraic
implication proved above. Conversely, the written argument does not assert
that a finite set of tested values alone proves an infinite formula: the
entrywise identities and their thresholds are the step that supplies the
tail.

\section{Discussion}

The three widths have different eventual periods: four for widths five and
seven, and fourteen for width six. This difference is visible in the finite
state graphs and in the residue correction for width six. The method scales
naturally as a fixed-width transfer construction, but the finite state space
grows rapidly with the width. No claim about width eight is made here.

The width-five and width-seven formulas have the same eventual increment
pattern but arise from different state spaces, while width six has a longer
period and a nontrivial residue correction. The period-four tails are visible
in the corresponding entrywise identities, whereas the period fourteen in
width six is an exact property of the certified matrix power, not an assumed
periodicity ansatz. The exceptional values are consequently treated as
finite-prefix data: they are not extrapolated from a recurrence whose base
threshold excludes them.

The contribution should be read as a fixed-width frontier result and a
reusable verification pattern. Earlier work settles smaller widths or
particular residue subfamilies, and nearby 2-domination work demonstrates
that tropical transfer methods are useful for cylindrical problems. The
present package adds complete all-circumference total-domination formulas for
the three frozen widths, together with a clean-room certificate protocol that
separates graph semantics, finite matrix identities, and tail induction. No
claim is made that the transfer architecture itself is new in general, or
that the method resolves arbitrary width.

A cyclic domination problem can be expressed as a closed-walk problem without
sacrificing the open-neighborhood condition. Once the finite weighted graph is
fixed, an entrywise min-plus identity controls every subsequent walk length at
and above its threshold; the finite prefix accounts for the transients and
exceptional circumferences.

\section{Conclusion}

We determined the exact total domination numbers of \(P_5\sq C_n\),
\(P_6\sq C_n\), and \(P_7\sq C_n\) for every \(n\geq3\). The proof combines a
weight-preserving finite-state bijection, min-plus path semantics, three exact
entrywise matrix identities, and finite-prefix verification. The
computer-assisted portion is specified by the finite proposition and the
reproducibility record, while the graph-theoretic reduction and recurrence
argument are given explicitly above.

\begin{thebibliography}{99}
\bibitem{Gravier2002}
S. Gravier, Total domination number of grid graphs, Discrete Appl. Math. 121
(2002) 119--128. \url{https://doi.org/10.1016/S0166-218X(01)00297-9}.

\bibitem{CrevalsOstergard2017}
S. Crevals, P. R. J. Ostergard, Total domination of grid graphs, J. Combin.
Math. Combin. Comput. 101 (2017) 175--192.

\bibitem{Eakawinrujee2022}
P. Eakawinrujee, Total and paired domination numbers of cylinders, Bull.
Malays. Math. Sci. Soc. 45 (2022) 3321--3334.
\url{https://doi.org/10.1007/s40840-022-01382-1}.

\bibitem{EakawinrujeeThesis2023}
P. Eakawinrujee, \emph{Total and paired domination numbers and $\gamma$-total
and $\gamma$-paired dominating graphs of some graphs}, Ph.D. dissertation,
Thammasat University, 2023.
\url{https://ethesisarchive.library.tu.ac.th/thesis/2023/TU_2023_6109320413_17206_27705.pdf}.

\bibitem{HuSohnChen2016}
F.-T. Hu, M. Y. Sohn, X.-g. Chen, Total and paired domination numbers of
$C_m$ bundles over a cycle $C_n$, J. Combin. Optim. 32 (2016) 608--625.
\url{https://doi.org/10.1007/s10878-015-9885-7}.

\bibitem{NandiParuiAdhikari2011}
M. Nandi, S. Parui, A. Adhikari, The domination numbers of cylindrical grid
graphs, Appl. Math. Comput. 217 (2011) 4879--4889.
\url{https://doi.org/10.1016/j.amc.2010.11.019}.

\bibitem{MartinezCastanoPuertas2024}
J. A. Martínez, A. B. Castaño-Fernández, M. L. Puertas, The 2-domination
number of cylindrical graphs, arXiv:2409.16703 (2024).
\url{https://arxiv.org/abs/2409.16703}.

\bibitem{GarzonMartinezMorenoPuertas2021}
E. M. Garzón, J. A. Martínez, J. J. Moreno, M. L. Puertas, On the
2-domination number of cylinders with small cycles, Fundam. Inform. 185
(2022) 185--199.
\url{https://arxiv.org/abs/2109.10549}.
\end{thebibliography}

\end{document}
